\section[Irreducibility of U(g)\^{}{G'}-modules]{Irreducibility of $\boldsymbol{\univ{g}^{G'}}$-modules}

In this section, we consider the irreducibility of $\univ{g}^{G'}$-modules in the branching problem.
We treat generalized Verma modules, cohomologically induced modules and discrete series representations.

\subsection{Generalized Verma module}

In this subsection, we consider the branching problem of generalized Verma modules.
We will give a criterion for the irreducibility of $\univ{g}^{G'}$-modules on $\Hom$ spaces.

Let $(\lie{g}, G')$ be a pair with semisimple $\lie{g}$.
Let $\lie{q} = \lie{l}\oplus \lie{u}$ be a weakly $\lie{g'}$-compatible parabolic subalgebra of $\lie{g}$
containing $\lie{q}(H)$ ($H \in \lie{g'}$).
See Definition~\ref{def:WeaklyCompatible}.
Set
\begin{align*}
	\lie{l'}\coloneq \lie{l}\cap \lie{g'}, \qquad \lie{u'}\coloneq \lie{u}\cap \lie{g'}, \qquad \lie{\overline{u}'}\coloneq \lie{\overline{u}}\cap \lie{g'}, \qquad \lie{q'}\coloneq \lie{l'}\oplus \lie{u'}.
\end{align*}
Then $\lie{q'}$ is a parabolic subalgebra of $\lie{g'}$.
Write $L'$ for the centralizer of $H$ in $G'$.

Fix a Cartan subalgebra $\lie{t}$ of $\lie{l}$ and write $W_{G}$ and $W_L$ for the Weyl groups.
Let $\rhog{\lie{u}}$ denote half the sum of roots in $\Delta(\lie{u}, \lie{t})$.
Similarly, we define $\lie{t'}$, $W_{G'}$, $W_{L'}$ and $\rhog{\lie{u'}}$ for $\lie{g'}$.
Fix a~set $\Delta^+(\lie{l'}, \lie{t'})$ of positive roots in $\Delta(\lie{l'}, \lie{t'})$
and $\rhog{\lie{l'}}$ denotes half the sum of the positive roots.
Set~${\rhog{\lie{g'}} \coloneq \rhog{\lie{l'}} + \rhog{\lie{u'}}}$.

In this subsection, using the results in Section~\ref{sect:Embedding}, we consider the irreducibility of the $\univ{g}^{G'}$-module
$\Hom_{\lie{g'}}\bigl(\ind{g'}{q'}(F'), \ind{g}{q}(F)|_{\lie{g'}}\bigr)$.
For a finite-dimensional irreducible $\lie{l}'$-module $F$, we set
\begin{align*}
	\fdual{F} \coloneq F^* \otimes \CC_{-2\rhog{\lie{u'}}}.
\end{align*}
Then the multiplicity of $\CC_{-2\rhog{\lie{u'}}}$ in $\fdual{F}\otimes F$ is $1$, and
\begin{align*}
	(-2\rhog{\lie{u'}} - \rhog{\lie{l'}} + \rhog{\lie{u'}}, \beta) = (-\rhog{\lie{l'}} - \rhog{\lie{u'}}, \beta) < 0, \qquad \forall \beta \in \Delta(\lie{u'}, \lie{t'}).
\end{align*}
Note that $-2\rhog{\lie{u'}} - \rhog{\lie{l'}}$ is the infinitesimal character of $\CC_{-2\rhog{\lie{u'}}}$
and $-\rhog{\lie{l'}} - \rhog{\lie{u'}} = -\rhog{\lie{g'}}$ is that of the $\lie{g'}$-module $\CC$.

\begin{Lemma}\label{lem:UniqueTrivialCharacter}
	Let $F$ be a finite-dimensional irreducible $\lie{l}'$-module and $F'$ an irreducible $\lie{l}'$-submodule of $\fdual{F}\otimes F$.
	If \smash{$\ind{g'}{q'}(F')$} has the infinitesimal character $[\rhog{\lie{g'}}]$, then $F' \simeq \CC_{-2\rhog{\lie{u'}}}$ holds.
\end{Lemma}

\begin{proof}
	Write $\lambda\in (\lie{t'})^*$ for the representative of the infinitesimal character of $F'$ such that~$\lambda$ is anti-dominant with respect to $\Delta^+(\lie{l'}, \lie{t'})$.
	Then $[\lambda + \rhog{\lie{u'}}]$ is the infinitesimal character of~\smash{$\ind{g'}{q'}(F')$}.
	By assumption, there exists $w \in W_{G'}$ such that $w(\lambda + \rhog{\lie{u'}}) = -\rhog{\lie{g'}}$.
	Since~$F'$ is a~submodule of $\fdual{F}\otimes F$, we have $\lambda(H) = (-2\rhog{\lie{u'}} - \rhog{\lie{l'}})(H)$, and hence
	\begin{align}
		(\lambda + \rhog{\lie{u'}})(H) = -\rhog{\lie{g'}}(H). \label{eqn:UniqueTrivialCharacter}
	\end{align}

	Let $\lie{b'}$ denote the Borel subalgebra of $\lie{g'}$ such that $-\rhog{\lie{g'}}$ is dominant with respect to $\lie{b'}$.
	Then $\lambda + \rhog{\lie{u'}}$ is dominant with respect to $w^{-1}\lie{b'}$.
	This implies that $\lambda + \rhog{\lie{u'}}$ is half the sum of roots in $\Delta\bigl(w^{-1}\lie{b'}, \lie{t'}\bigr)$.
	By \eqref{eqn:UniqueTrivialCharacter}, we have $\lie{\overline{u}'} \subset w^{-1}\lie{b'}$.
	Since $\lambda$ is dominant with respect to $-\Delta^+(\lie{l'}, \lie{t'})$, this implies $w^{-1}\lie{b'} = \lie{b'}$
	and hence $w = e$.
	Therefore, we obtain $\lambda = -2\rhog{\lie{u'}} - \rhog{\lie{l'}}$ and $F'\simeq \CC_{-2\rhog{\lie{u'}}}$.
\end{proof}

By Lemma~\ref{lem:UniqueTrivialCharacter} and the proof of Lemma~\ref{lem:MinimalTypeOne}, we obtain the following.

\begin{Lemma}\label{lem:ExistenceInvariantVector}
	Let $F$ be a finite-dimensional irreducible $\lie{l}'$-module and $M$ a non-zero quotient of~${\ind{g'}{q'}\bigl(\fdual{F}\bigr)\otimes \ind{g'}{q'}(F)}$.
	\begin{enumerate}\itemsep=0pt
		\item[$1.$] \smash{$P_{-\rhog{\lie{l'}} - \rhog{\lie{u'}}}\bigl(M|_{\Delta(\lie{g'})}\bigr) \simeq \ind{g'}{q'}\bigl(\CC_{-2\rhog{\lie{u'}}}\bigr)$}.
		\item[$2.$] \smash{$\Dzuck{\Delta(G')}{\Delta(L')}{S}(M)^{\Delta(G')} \simeq \CC$}.
		\item[$3.$] \smash{$\Dzuck{\Delta(G')}{\Delta(L')}{i}(M)^{\Delta(G')} = 0$} for any $i \neq S$.
	\end{enumerate}
\end{Lemma}

By Lemma~\ref{lem:ExistenceInvariantVector}, the functor \smash{$\Dzuck{\Delta(G')}{\Delta(L')}{S}\bigl(\ind{g'}{q'}\bigl(\fdual{F}\bigr)\otimes (\cdot)\bigr)^{\Delta(G')}$} is regarded as an analogue of the functor \smash{$\Hom_{\lie{g'}}\bigl(\ind{g'}{q'}(F), \cdot\bigr)$}.
We shall compare the two functors.

Let $F'$ a finite-dimensional irreducible $\lie{l'}$-module and $M \in \BGGcat{g}{q}$.
Suppose that the $\lie{l}'$-module $M \otimes (F')^*$ lifts to an $L'$-module.
This implies that the $\lie{l}'$-module \smash{$\ind{g'}{q'}\bigl(\fdual{F'}\bigr)\otimes M$} lifts to an~\mbox{$L'$-module}.
Remark that this condition follows from \smash{$\Hom_{\lie{g'}}\bigl(\ind{g'}{q'}(F'), M\bigr) \neq 0$} if $M$ is indecomposable.

There is a canonical $\lie{g'}$- and $\univ{g}^{G'}$-homomorphism
\begin{align*}
\Phi \colon\ \ind{g'}{q'}(F') \otimes \Hom_{\lie{g'}}\bigl(\ind{g'}{q'}(F'), M\bigr)
\rightarrow M
\end{align*}
defined by $\Phi(v \otimes \varphi) = \varphi(v)$.
Let $M'$ be a non-zero quotient of \smash{$\ind{g'}{q'}\bigl(\fdual{F'}\bigr)$}.
Applying Lem\-ma~\ref{lem:ExistenceInvariantVector}\,(2) to \smash{$M'\otimes \ind{g'}{q'}(F')$}, we have isomorphisms
\begin{align}
	&\Hom_{\lie{g'}}\bigl(\ind{g'}{q'}(F'), M\bigr) \simeq \CC\otimes \Hom_{\lie{g'}}\bigl(\ind{g'}{q'}(F'), M\bigr) \nonumber \\
	&\qquad{}\simeq \Dzuck{\Delta(G')}{\Delta(L')}{S}\bigl(M'\otimes \ind{g'}{q'}(F')\bigr)^{\Delta(G')}\otimes \Hom_{\lie{g'}}\bigl(\ind{g'}{q'}(F'), M\bigr) \nonumber \\
	&\qquad{}\simeq \Dzuck{\Delta(G')}{\Delta(L')}{S}\bigl(M'\otimes \ind{g'}{q'}(F') \otimes \Hom_{\lie{g'}}\bigl(\ind{g'}{q'}(F'), M\bigr)\bigr)^{\Delta(G')} \label{eqn:IsomHomZuck}
\end{align}
of $\univ{g}^{G'}$-modules.
Compositing this isomorphism with \smash{$\Dzuck{\Delta(G')}{\Delta(L')}{S}(\id\otimes \Phi)^{\Delta(G')}$} gives a~$\univ{g}^{G'}$-homomorphism
\begin{align*}
	\Phi_{M'} \colon\ \Hom_{\lie{g'}}\bigl(\ind{g'}{q'}(F'), M\bigr)\rightarrow \Dzuck{\Delta(G')}{\Delta(L')}{S}(M'\otimes M)^{\Delta(G')}.
\end{align*}
If $\Phi_{M'}$ is injective, we can study the space $\Hom_{\lie{g'}}\bigl(\ind{g'}{q'}(F'), M\bigr)$ using the Zuckerman derived functor module.
We shall consider the injectivity of $\Phi_{M'}$.

\begin{Lemma}\label{lem:VanishingInvariantKernel}
	\smash{$\Dzuck{\Delta(G')}{\Delta(L')}{i}(M'\otimes \Ker(\Phi))^{\Delta(G')} = 0$} holds for any $i$.
\end{Lemma}

\begin{proof}
	Let $N$ denote the unique maximal submodule of $\ind{g'}{q'}(F')$.
	By Lemma~\ref{lem:ExistenceInvariantVector}, it is enough to show
	\begin{align}
	\Ker(\Phi) \subset N \otimes \Hom_{\lie{g'}}\bigl(\ind{g'}{q'}(F'), M\bigr). \label{eqn:VanishingInvariantKernel}
	\end{align}

	Let $\mu$ denote the highest weight of $\ind{g'}{q'}(F')$.
	Assume that \eqref{eqn:VanishingInvariantKernel} does not hold.
	Set $X \coloneq \Hom_{\lie{g'}}\bigl(\ind{g'}{q'}(F'), M\bigr)$.
	Then we have
	\begin{align*}
		0\neq \Ker(\Phi)/(\Ker(\Phi) \cap (N \otimes X)) \subset \bigl(\ind{g'}{q'}(F') / N\bigr) \otimes X.
	\end{align*}
	Since \smash{$\bigl(\ind{g'}{q'}(F')/N\bigr) \otimes X$} is isomorphic to a direct sum of some copies of the irreducible $\lie{g'}$-mod\-ule $\ind{g'}{q'}(F')/N$ as a $\lie{g'}$-module, there is a~weight vector $v \in \Ker(\Phi)$ with the weight $\mu$.
	Then $v$ is a~$\lie{b'}$-eigenvector.

	Write $v= \sum_{i=0}^r v_i\otimes \varphi_i$, where \smash{$\set{v_i}_{i=0}^r \subset \ind{g'}{q'}(F')$}
	and $\set{\varphi_i}_{i=0}^r \subset X$
	such that~$\set{v_i}_{i=0}^r$ and~$\set{\varphi_i}_{i=0}^r$ are linearly independent.
	Then each $v_i$ is a highest weight vector of \smash{$\ind{g'}{q'}(F')$}.
	This implies $r=0$.
	By $v \in \Ker(\Phi)$, we have $\varphi_0(v_0)=0$.
	Since $v_0$ is a cyclic vector of \smash{$\ind{g'}{q'}(F')$}, we have $\varphi_0 = 0$ and hence $v = 0$.
	This contradicts to $v \neq 0$.
	Therefore, we have shown \eqref{eqn:VanishingInvariantKernel}.
\end{proof}

Consider the following two exact sequences:
\begin{align*}
	&0 \rightarrow \Ker(\Phi)\rightarrow \ind{g'}{q'}(F') \otimes \Hom_{\lie{g'}}\bigl(\ind{g'}{q'}(F'), M\bigr) \xrightarrow{\Phi} \Im(\Phi) \rightarrow 0, \\
	&0 \rightarrow \Im(\Phi)\rightarrow M \rightarrow \Coker(\Phi) \rightarrow 0.
\end{align*}
Applying the functor \smash{$\Dzuck{\Delta(G')}{\Delta(L')}{i}(M'\otimes (\cdot))^{\Delta(G')}$}, and by Lemma~\ref{lem:VanishingInvariantKernel} and \eqref{eqn:IsomHomZuck}, we obtain
\begin{align*}
	&\Dzuck{\Delta(G')}{\Delta(L')}{S}(M'\otimes \Im(\Phi))^{\Delta(G')}\simeq \Hom_{\lie{g'}}\bigl(\ind{g'}{q'}(F'), M\bigr), \\
	&\Dzuck{\Delta(G')}{\Delta(L')}{S - 1}(M'\otimes \Coker(\Phi))^{\Delta(G')}
 \rightarrow \Dzuck{\Delta(G')}{\Delta(L')}{S}(M'\otimes \Im(\Phi))^{\Delta(G')} \\
	&\hphantom{\Dzuck{\Delta(G')}{\Delta(L')}{S - 1}(M'\otimes \Coker(\Phi))^{\Delta(G')}}{}
 \xrightarrow{\Phi_{M'}} \Dzuck{\Delta(G')}{\Delta(L')}{S}(M'\otimes M)^{\Delta(G')} \qquad (\text{exact})
\end{align*}
as $\univ{g}^{G'}$-modules.
Therefore, we have shown the following criterion.

\begin{Theorem}\label{thm:InjectivePhi}
	$\Phi_{M'}$ is injective if and only if \smash{$\Dzuck{\Delta(G')}{\Delta(L')}{S - 1}(M'\otimes \Coker(\Phi))^{\Delta(G')} = 0$}.
\end{Theorem}

\begin{Remark}\label{rmk:InjectivePhiStandard}
	If $(M'\otimes \Coker(\Phi))|_{\Delta(\lie{g'})}$ has an exhaustive filtration whose associated graded module has a standard filtration, we have \smash{$\Dzuck{\Delta(G')}{\Delta(L')}{S - 1}(M'\otimes \Coker(\Phi))^{\Delta(G')} = 0$} by Proposition~\ref{prop:StandardFiltrationVanishing}.
\end{Remark}

